\documentclass[12pt,a4paper]{article}
\usepackage[spanish]{babel}
\usepackage[utf8]{inputenc}
\usepackage{amsmath,amssymb,amsthm}
\usepackage{graphicx}
\usepackage{tikz}
\usepackage[margin=2.5cm]{geometry}
\usepackage{enumitem}
\usepackage{multicol}
\usepackage{caption}
\usepackage{hyperref}

% Para graficar funciones trigonométricas inversas en radianes
\usetikzlibrary{arrows.meta}

\newtheorem{ejemplo}{Ejemplo}
\newtheorem*{sol}{Solución}

\title{\textbf{Guía de Ecuaciones Trigonométricas Cuadráticas\\ 
y Funciones Trigonométricas Inversas}}
\author{}
\date{}

\begin{document}
\maketitle

\section{Introducción}
En esta guía aprenderás a resolver ecuaciones trigonométricas que pueden transformarse en ecuaciones cuadráticas mediante un cambio de variable. Además, repasaremos las funciones trigonométricas inversas, sus dominios, recorridos y gráficas, herramientas indispensables para expresar soluciones no exactas en el intervalo $[0,2\pi)$.

\section{Ecuaciones trigonométricas reducibles a cuadráticas}
El método consiste en:
\begin{enumerate}
    \item Identificar una función trigonométrica que aparece repetida, por ejemplo $\sin x$, $\cos x$, $\tan x$, etc.
    \item Realizar el cambio $u = \text{función trigonométrica}$.
    \item Resolver la ecuación cuadrática resultante en $u$.
    \item Volver a la variable original y resolver las ecuaciones trigonométricas elementales.
    \item Expresar todas las soluciones en el intervalo $[0,2\pi)$, usando funciones trigonométricas inversas si los valores no son ángulos especiales.
\end{enumerate}

\begin{ejemplo}
Resolver $2\sin^2 x - 3\sin x + 1 = 0$ en $[0,2\pi)$.
\end{ejemplo}
\begin{sol}
Sea $u = \sin x$. Entonces la ecuación se transforma en
\[
2u^2 - 3u + 1 = 0,
\]
que factoriza como $(2u-1)(u-1)=0$. Luego $u = \frac{1}{2}$ o $u = 1$.
Volviendo a la variable original:
\begin{itemize}
    \item $\sin x = \frac{1}{2} \;\Rightarrow\; x = \frac{\pi}{6},\; \frac{5\pi}{6}$ (en $[0,2\pi)$).
    \item $\sin x = 1 \;\Rightarrow\; x = \frac{\pi}{2}$.
\end{itemize}
Conjunto solución: $\left\{\frac{\pi}{6},\; \frac{\pi}{2},\; \frac{5\pi}{6}\right\}$.
\end{sol}

\section{Funciones trigonométricas inversas}
Para manejar soluciones no exactas necesitamos las funciones inversas $\arcsin$, $\arccos$ y $\arctan$. Se definen restringiendo el dominio de las funciones originales para que sean biyectivas.

\subsection{Arcoseno, $y = \arcsin x$}
\begin{minipage}{0.45\textwidth}
    \begin{itemize}
        \item \textbf{Dominio:} $[-1,1]$
        \item \textbf{Recorrido:} $\left[-\frac{\pi}{2},\frac{\pi}{2}\right]$
        \item \textbf{Definición:} $\arcsin x = y \iff \sin y = x$,\, $y\in\left[-\frac{\pi}{2},\frac{\pi}{2}\right]$.
    \end{itemize}
\end{minipage}\hfill
\begin{minipage}{0.5\textwidth}
    \centering
    \begin{tikzpicture}[scale=0.8]
        \draw[->] (-1.5,0) -- (1.5,0) node[right] {$x$};
        \draw[->] (0,-2) -- (0,2) node[above] {$y$};
        \draw[domain=-1:1,smooth,variable=\x,blue,thick] plot ({\x},{rad(asin(\x))});
        \node at (1.2, {pi/2}) [right] {$\frac{\pi}{2}$};
        \node at (-1.2, {-pi/2}) [left] {$-\frac{\pi}{2}$};
        \foreach \x in {-1,1} \draw (\x,0.1) -- (\x,-0.1) node[below] {$\x$};
        \node at (0,-2.5) {$y = \arcsin x$};
    \end{tikzpicture}
\end{minipage}

\subsection{Arcocoseno, $y = \arccos x$}
\begin{minipage}{0.45\textwidth}
    \begin{itemize}
        \item \textbf{Dominio:} $[-1,1]$
        \item \textbf{Recorrido:} $[0,\pi]$
        \item \textbf{Definición:} $\arccos x = y \iff \cos y = x$,\, $y\in[0,\pi]$.
    \end{itemize}
\end{minipage}\hfill
\begin{minipage}{0.5\textwidth}
    \centering
    \begin{tikzpicture}[scale=0.8]
        \draw[->] (-1.5,0) -- (1.5,0) node[right] {$x$};
        \draw[->] (0,-0.5) -- (0,4) node[above] {$y$};
        \draw[domain=-1:1,smooth,variable=\x,blue,thick] plot ({\x},{rad(acos(\x))});
        \node at (1.2, 0) [right] {$0$};
        \node at (-1.2, {pi}) [left] {$\pi$};
        \foreach \x in {-1,1} \draw (\x,0.1) -- (\x,-0.1) node[below] {$\x$};
        \node at (0,-0.8) {$y = \arccos x$};
    \end{tikzpicture}
\end{minipage}

\subsection{Arcotangente, $y = \arctan x$}
\begin{minipage}{0.45\textwidth}
    \begin{itemize}
        \item \textbf{Dominio:} $\mathbb{R}$
        \item \textbf{Recorrido:} $\left(-\frac{\pi}{2},\frac{\pi}{2}\right)$
        \item \textbf{Definición:} $\arctan x = y \iff \tan y = x$,\, $y\in\left(-\frac{\pi}{2},\frac{\pi}{2}\right)$.
    \end{itemize}
\end{minipage}\hfill
\begin{minipage}{0.5\textwidth}
    \centering
    \begin{tikzpicture}[scale=0.8]
        \draw[->] (-3.5,0) -- (3.5,0) node[right] {$x$};
        \draw[->] (0,-2) -- (0,2) node[above] {$y$};
        \draw[domain=-3:3,smooth,variable=\x,blue,thick] plot ({\x},{rad(atan(\x))});
        \draw[dashed] (-3,{pi/2}) -- (3,{pi/2});
        \draw[dashed] (-3,{-pi/2}) -- (3,{-pi/2});
        \node at (0.3, {pi/2}) [left] {$\frac{\pi}{2}$};
        \node at (0.3, {-pi/2}) [left] {$-\frac{\pi}{2}$};
        \node at (0,-2.5) {$y = \arctan x$};
    \end{tikzpicture}
\end{minipage}

\subsection{Fórmulas generales para soluciones en $[0,2\pi)$}
Dado un valor $a$ dentro del dominio adecuado:
\begin{align*}
\sin x &= a \quad \Rightarrow \quad x = 
\begin{cases}
\arcsin a + 2k\pi,\\
\pi - \arcsin a + 2k\pi,
\end{cases} k\in\mathbb{Z}.\\
\cos x &= a \quad \Rightarrow \quad x = 
\pm \arccos a + 2k\pi,\; k\in\mathbb{Z}.\\
\tan x &= a \quad \Rightarrow \quad x = 
\arctan a + k\pi,\; k\in\mathbb{Z}.
\end{align*}
Para obtener las soluciones en $[0,2\pi)$, se toma $k=0,1$ o el valor que corresponda, descartando las que caigan fuera del intervalo.

\newpage
\section{Ejercicios propuestos}
Resuelve cada ecuación en el intervalo $[0,2\pi)$. Cuando la solución no corresponda a un ángulo especial, exprésala usando funciones trigonométricas inversas.

\begin{enumerate}[leftmargin=*,label=\arabic*.]

\item $2\cos^2 x - \cos x - 1 = 0$

\item $3\tan^2 x - 4\tan x + 1 = 0$

\item $\sin^2 x + \sin x - 2 = 0$

\item $2\cos^2 x - 5\cos x + 2 = 0$

\item $4\sin^2 x - 4\sin x + 1 = 0$

\item $\sec^2 x - 3\sec x + 2 = 0$

\item $2\csc^2 x - 3\csc x - 2 = 0$

\item $3\sin^2 x + \sin x - 2 = 0$

\item $1 - \cos^2 x - \cos x = 0$ \quad (Sugerencia: usa $\sin^2 x = 1-\cos^2 x$)

\item $2\sin^2 x + 5\sin x - 3 = 0$

\item $\tan^2 x + 2\tan x - 3 = 0$

\item $\cos^2 x - \sin x = 0$

\item $3\cos^2 x + 2\cos x - 1 = 0$

\item $2\sec^2 x + 5\sec x + 2 = 0$

\item $2\csc^2 x - 7\csc x + 3 = 0$

\item $3\tan^2 x + \tan x - 2 = 0$

\item $\cos^2 x - 2\cos x - 2 = 0$

\item $5\sin^2 x - 4\sin x - 1 = 0$

\item $3\tan^2 x - 1 = 0$

\item $2\sin^2 x - \sin x - 1 = 0$
\end{enumerate}

\newpage
\section*{Soluciones}
A continuación se presentan las soluciones finales en $[0,2\pi)$. Para los ángulos no exactos se ha empleado la notación de funciones inversas.

\begin{multicols}{2}
\begin{enumerate}[leftmargin=*,label=\arabic*.]
    \item $\displaystyle 0,\; \frac{2\pi}{3},\; \frac{4\pi}{3}$
    \item $\displaystyle \frac{\pi}{4},\; \frac{5\pi}{4},\; \arctan\!\left(\frac{1}{3}\right),\; \pi + \arctan\!\left(\frac{1}{3}\right)$
    \item $\displaystyle \frac{\pi}{2}$
    \item $\displaystyle \frac{\pi}{3},\; \frac{5\pi}{3}$
    \item $\displaystyle \frac{\pi}{6},\; \frac{5\pi}{6}$
    \item $\displaystyle 0,\; \frac{\pi}{3},\; \frac{5\pi}{3}$
    \item $\displaystyle \frac{\pi}{6},\; \frac{5\pi}{6}$
    \item $\displaystyle \frac{3\pi}{2},\; \arcsin\!\left(\frac{2}{3}\right),\; \pi - \arcsin\!\left(\frac{2}{3}\right)$
    \item $\displaystyle \arccos\!\left(\frac{\sqrt{5}-1}{2}\right),\; 2\pi - \arccos\!\left(\frac{\sqrt{5}-1}{2}\right)$
    \item $\displaystyle \frac{\pi}{6},\; \frac{5\pi}{6}$
    \item $\displaystyle \frac{\pi}{4},\; \frac{5\pi}{4},\; \pi - \arctan 3,\; 2\pi - \arctan 3$
    \item $\displaystyle \arcsin\!\left(\frac{\sqrt{5}-1}{2}\right),\; \pi - \arcsin\!\left(\frac{\sqrt{5}-1}{2}\right)$
    \item $\displaystyle \pi,\; \arccos\!\left(\frac{1}{3}\right),\; 2\pi - \arccos\!\left(\frac{1}{3}\right)$
    \item $\displaystyle \frac{2\pi}{3},\; \frac{4\pi}{3}$
    \item $\displaystyle \arcsin\!\left(\frac{1}{3}\right),\; \pi - \arcsin\!\left(\frac{1}{3}\right)$
    \item $\displaystyle \frac{3\pi}{4},\; \frac{7\pi}{4},\; \arctan\!\left(\frac{2}{3}\right),\; \pi + \arctan\!\left(\frac{2}{3}\right)$
    \item $\displaystyle \arccos(1-\sqrt{3}),\; 2\pi - \arccos(1-\sqrt{3})$
    \item $\displaystyle \frac{\pi}{2},\; \pi + \arcsin\!\left(\frac{1}{5}\right),\; 2\pi - \arcsin\!\left(\frac{1}{5}\right)$
    \item $\displaystyle \frac{\pi}{6},\; \frac{5\pi}{6},\; \frac{7\pi}{6},\; \frac{11\pi}{6}$
    \item $\displaystyle \frac{\pi}{2},\; \frac{7\pi}{6},\; \frac{11\pi}{6}$
\end{enumerate}
\end{multicols}

\vspace{1cm}
{\small Nota: En los ejercicios 12, 13 y 17 se usó la identidad pitagórica para convertir la ecuación a una cuadrática en una sola función. En todos los casos, las soluciones se ajustaron al intervalo $[0,2\pi)$ descartando ángulos fuera de rango.}

\end{document}